\documentclass[12pt,a4paper]{article}
\usepackage[UTF8]{ctex}
\usepackage{geometry}
\usepackage{hyperref}
\usepackage{listings}
\usepackage{xcolor}
\usepackage{booktabs}
\usepackage{longtable}
\usepackage{enumitem}
\usepackage{fancyhdr}

\geometry{left=2.5cm,right=2.5cm,top=3cm,bottom=3cm}

% 代码样式
\lstset{
  basicstyle=\ttfamily\small,
  breaklines=true,
  frame=single,
  backgroundcolor=\color{gray!10},
  keywordstyle=\color{blue},
  stringstyle=\color{red},
  commentstyle=\color{green!60!black},
  numbers=left,
  numberstyle=\tiny\color{gray},
}

% 定义 JSON 语言
\lstdefinelanguage{json}{
  basicstyle=\ttfamily\small,
  stringstyle=\color{red},
  commentstyle=\color{green!60!black},
  keywordstyle=\color{blue},
  numbers=left,
  numberstyle=\tiny\color{gray},
  breaklines=true,
  frame=single,
  backgroundcolor=\color{gray!10},
}

% 定义 JavaScript 语言
\lstdefinelanguage{javascript}{
  basicstyle=\ttfamily\small,
  keywordstyle=\color{blue}\bfseries,
  stringstyle=\color{red},
  commentstyle=\color{green!60!black},
  numbers=left,
  numberstyle=\tiny\color{gray},
  breaklines=true,
  frame=single,
  backgroundcolor=\color{gray!10},
  morekeywords={const, let, var, function, async, await, fetch, JSON, console, log},
}

% 页眉页脚
\pagestyle{fancy}
\setlength{\headheight}{14.5pt}
\fancyhf{}
\fancyhead[L]{智能公交安全助理 - API接口文档}
\fancyhead[R]{\today}
\fancyfoot[C]{\thepage}

\title{智能公交安全助理\\API接口文档}
\author{Agent Service API}
\date{\today}

\begin{document}

\maketitle
\tableofcontents
\newpage

\section{概述}

\subsection{简介}

智能公交安全助理（Agent Service）提供基于大语言模型的智能对话服务，支持生成安全评估报告、智能咨询等功能。本文档描述了Agent服务的REST API接口规范。

\subsection{基础信息}

\begin{itemize}
  \item \textbf{基础路径}: \texttt{/api/agent}
  \item \textbf{协议}: HTTP/HTTPS
  \item \textbf{数据格式}: JSON
  \item \textbf{字符编码}: UTF-8
\end{itemize}

\subsection{版本信息}

\begin{itemize}
  \item \textbf{API版本}: v1
  \item \textbf{文档版本}: 1.0.0
  \item \textbf{最后更新}: \today
\end{itemize}

\section{通用规范}

\subsection{请求格式}

所有请求应包含以下HTTP头：

\begin{lstlisting}[language=bash]
Content-Type: application/json
Accept: application/json
\end{lstlisting}

\subsection{响应格式}

所有响应均为JSON格式，包含以下字段：

\begin{itemize}
  \item \textbf{成功响应}: 直接返回数据对象或数组
  \item \textbf{错误响应}: 包含 \texttt{error} 字段，描述错误信息
\end{itemize}

\subsection{HTTP状态码}

\begin{table}[h]
  \centering
  \begin{tabular}{ll}
    \toprule
    状态码 & 说明 \\
    \midrule
    200 & 请求成功 \\
    201 & 创建成功 \\
    400 & 请求参数错误 \\
    404 & 资源不存在 \\
    500 & 服务器内部错误 \\
    \bottomrule
  \end{tabular}
  \caption{HTTP状态码说明}
\end{table}

\subsection{错误响应格式}

\begin{lstlisting}[language=json]
{
  "error": "错误描述信息"
}
\end{lstlisting}

\section{接口列表}

\subsection{系统接口}

\subsubsection{健康检查}

检查服务运行状态。

\textbf{请求}

\begin{lstlisting}[language=bash]
GET /api/agent/health
\end{lstlisting}

\textbf{响应}

\begin{lstlisting}[language=json]
{
  "status": "ok",
  "timestamp": "2024-01-15T10:30:00.000Z"
}
\end{lstlisting}

\textbf{状态码}: 200

\subsection{对话接口}

\textbf{重要说明}：Agent服务采用智能路由机制，会根据用户消息的意图自动识别并分发到相应的工具进行处理。调用端无需关心具体的工具列表，只需发送消息即可。

\subsubsection{获取会话列表}

获取所有会话列表，按更新时间倒序排列。

\textbf{请求}

\begin{lstlisting}[language=bash]
GET /api/agent/sessions
\end{lstlisting}

\textbf{响应}

\begin{lstlisting}[language=json]
[
  {
    "id": "session_123456",
    "title": "新会话",
    "preview": "请生成张三的安全报告...",
    "updatedAt": "2024-01-15T10:30:00.000Z"
  }
]
\end{lstlisting}

\textbf{字段说明}

\begin{table}[h]
  \centering
  \begin{tabular}{lll}
    \toprule
    字段 & 类型 & 说明 \\
    \midrule
    id & String & 会话ID \\
    title & String & 会话标题 \\
    preview & String & 会话预览（最后一条消息的摘要） \\
    updatedAt & String & 更新时间（ISO 8601格式） \\
    \bottomrule
  \end{tabular}
  \caption{会话列表响应字段}
\end{table}

\textbf{状态码}: 200

\subsubsection{创建会话}

创建一个新的会话。

\textbf{请求}

\begin{lstlisting}[language=bash]
POST /api/agent/sessions
Content-Type: application/json

{
  "title": "新会话"  // 可选，默认为"新会话"
}
\end{lstlisting}

\textbf{请求参数}

\begin{table}[h]
  \centering
  \begin{tabular}{llll}
    \toprule
    参数 & 类型 & 必填 & 说明 \\
    \midrule
    title & String & 否 & 会话标题，默认为"新会话" \\
    \bottomrule
  \end{tabular}
  \caption{创建会话请求参数}
\end{table}

\textbf{响应}

\begin{lstlisting}[language=json]
{
  "id": "session_123456",
  "title": "新会话",
  "preview": "",
  "updatedAt": "2024-01-15T10:30:00.000Z"
}
\end{lstlisting}

\textbf{状态码}: 201

\subsubsection{获取会话详情}

获取指定会话的详细信息，包括所有消息记录。

\textbf{请求}

\begin{lstlisting}[language=bash]
GET /api/agent/sessions/{sessionId}
\end{lstlisting}

\textbf{路径参数}

\begin{table}[h]
  \centering
  \begin{tabular}{lll}
    \toprule
    参数 & 类型 & 说明 \\
    \midrule
    sessionId & String & 会话ID \\
    \bottomrule
  \end{tabular}
  \caption{获取会话详情路径参数}
\end{table}

\textbf{响应}

\begin{lstlisting}[language=json]
{
  "id": "session_123456",
  "title": "新会话",
  "preview": "请生成张三的安全报告...",
  "updatedAt": "2024-01-15T10:30:00.000Z",
  "messages": [
    {
      "id": "message_001",
      "role": "user",
      "content": "请生成张三的安全报告",
      "createdAt": "2024-01-15T10:25:00.000Z",
      "status": "complete"
    },
    {
      "id": "message_002",
      "role": "assistant",
      "content": "正在为您生成张三的安全报告...",
      "createdAt": "2024-01-15T10:25:05.000Z",
      "status": "complete",
      "metadata": {
        "tool": "generate_driver_report",
        "args": {
          "driver_name": "张三"
        }
      }
    }
  ]
}
\end{lstlisting}

\textbf{字段说明}

\begin{table}[h]
  \centering
  \begin{tabular}{lll}
    \toprule
    字段 & 类型 & 说明 \\
    \midrule
    id & String & 会话ID \\
    title & String & 会话标题 \\
    preview & String & 会话预览 \\
    updatedAt & String & 更新时间 \\
    messages & Array & 消息列表 \\
    messages[].id & String & 消息ID \\
    messages[].role & String & 消息角色（user/assistant/system/tool） \\
    messages[].content & String & 消息内容 \\
    messages[].createdAt & String & 创建时间 \\
    messages[].status & String & 状态（complete/error） \\
    messages[].metadata & Object & 元数据（可选） \\
    \bottomrule
  \end{tabular}
  \caption{会话详情响应字段}
\end{table}

\textbf{状态码}: 200

\textbf{错误响应}: 404（会话不存在）

\subsubsection{删除会话}

删除指定的会话及其所有消息记录。

\textbf{请求}

\begin{lstlisting}[language=bash]
DELETE /api/agent/sessions/{sessionId}
\end{lstlisting}

\textbf{路径参数}

\begin{table}[h]
  \centering
  \begin{tabular}{lll}
    \toprule
    参数 & 类型 & 说明 \\
    \midrule
    sessionId & String & 会话ID \\
    \bottomrule
  \end{tabular}
  \caption{删除会话路径参数}
\end{table}

\textbf{响应}

\begin{lstlisting}[language=json]
{
  "success": true
}
\end{lstlisting}

\textbf{状态码}: 200

\textbf{错误响应}: 404（会话不存在）

\subsection{对话接口}

\subsubsection{发送消息（非流式）}

向Agent发送消息，并获取AI助手的回复。此接口会等待完整响应后返回。

\textbf{请求}

\begin{lstlisting}[language=bash]
POST /api/agent/chat
Content-Type: application/json

{
  "content": "请生成张三的安全报告",
  "messages": [
    {
      "role": "user",
      "content": "你好"
    },
    {
      "role": "assistant",
      "content": "您好，我是智能公交安全助理，有什么可以帮助您的吗？"
    }
  ]
}
\end{lstlisting}

\textbf{请求参数}

\begin{table}[h]
  \centering
  \begin{tabular}{llll}
    \toprule
    参数 & 类型 & 必填 & 说明 \\
    \midrule
    content & String & 是 & 当前用户消息内容 \\
    messages & Array & 否 & 历史消息列表，用于构建上下文。每个消息包含role（user/assistant）和content字段 \\
    \bottomrule
  \end{tabular}
  \caption{发送消息请求参数}
\end{table}

\textbf{响应}

\begin{lstlisting}[language=json]
{
  "id": "message_002",
  "role": "assistant",
  "content": "正在为您生成张三的安全报告...",
  "createdAt": "2024-01-15T10:25:05.000Z",
  "status": "complete",
  "metadata": {
    "tool": "generate_driver_report",
    "args": {
      "driver_name": "张三"
    },
    "iterations": 2
  }
}
\end{lstlisting}

\textbf{字段说明}

\begin{table}[h]
  \centering
  \begin{tabular}{lll}
    \toprule
    字段 & 类型 & 说明 \\
    \midrule
    id & String & 消息ID \\
    role & String & 消息角色（assistant） \\
    content & String & AI回复内容 \\
    createdAt & String & 创建时间 \\
    status & String & 状态（complete/error） \\
    metadata & Object & 元数据（可选） \\
    metadata.tool & String & 使用的工具名称 \\
    metadata.args & Object & 工具参数 \\
    metadata.iterations & Number & 工具调用迭代次数 \\
    \bottomrule
  \end{tabular}
  \caption{发送消息响应字段}
\end{table}

\textbf{状态码}: 200

\textbf{错误响应}

\begin{itemize}
  \item 400: 缺少 sessionId 或 content
  \item 404: 会话不存在
\end{itemize}

\subsubsection{发送消息（流式）}

向指定会话发送消息，以Server-Sent Events (SSE)流式返回AI助手的回复。适用于需要实时显示回复内容的场景。

\textbf{请求}

\begin{lstlisting}[language=bash]
POST /api/agent/chat/stream
Content-Type: application/json

{
  "sessionId": "session_123456",
  "content": "请分析一下最近的风险趋势"
}
\end{lstlisting}

\textbf{请求参数}

\begin{table}[h]
  \centering
  \begin{tabular}{llll}
    \toprule
    参数 & 类型 & 必填 & 说明 \\
    \midrule
    sessionId & String & 是 & 会话ID \\
    content & String & 是 & 用户消息内容 \\
    \bottomrule
  \end{tabular}
  \caption{流式发送消息请求参数}
\end{table}

\textbf{响应格式}

响应为Server-Sent Events (SSE)流，包含以下事件类型：

\begin{enumerate}
  \item \textbf{start}: 开始处理
  \item \textbf{chunk}: 内容片段（可能多次）
  \item \textbf{complete}: 完成处理
  \item \textbf{error}: 错误信息（如果发生错误）
\end{enumerate}

\textbf{事件格式}

\begin{lstlisting}[language=bash]
# 开始事件
data: {"type":"start","messageId":"message_002"}

# 内容片段事件（可能多次）
data: {"type":"chunk","content":"正在"}

data: {"type":"chunk","content":"为您"}

data: {"type":"chunk","content":"生成"}

# 完成事件
data: {"type":"complete","messageId":"message_002","metadata":{...}}

# 错误事件（如果发生）
data: {"type":"error","error":"错误信息"}
\end{lstlisting}

\textbf{事件字段说明}

\begin{table}[h]
  \centering
  \begin{tabular}{lll}
    \toprule
    字段 & 类型 & 说明 \\
    \midrule
    type & String & 事件类型（start/chunk/complete/error） \\
    messageId & String & 消息ID（start/complete事件） \\
    content & String & 内容片段（chunk事件） \\
    metadata & Object & 元数据（complete事件） \\
    error & String & 错误信息（error事件） \\
    \bottomrule
  \end{tabular}
  \caption{SSE事件字段说明}
\end{table}

\textbf{HTTP头}

\begin{lstlisting}[language=bash]
Content-Type: text/event-stream
Cache-Control: no-cache
Connection: keep-alive
\end{lstlisting}

\textbf{状态码}: 200

\textbf{错误响应}: 404（会话不存在）

\textbf{客户端示例（JavaScript）}

\begin{lstlisting}[language=javascript]
const response = await fetch('/api/agent/chat/stream', {
  method: 'POST',
  headers: { 'Content-Type': 'application/json' },
  body: JSON.stringify({
    sessionId: 'session_123456',
    content: '请分析一下最近的风险趋势'
  })
});

const reader = response.body.getReader();
const decoder = new TextDecoder();

while (true) {
  const { done, value } = await reader.read();
  if (done) break;
  
  const chunk = decoder.decode(value);
  const lines = chunk.split('\n');
  
  for (const line of lines) {
    if (line.startsWith('data: ')) {
      const data = JSON.parse(line.slice(6));
      
      switch (data.type) {
        case 'start':
          console.log('开始处理:', data.messageId);
          break;
        case 'chunk':
          console.log('内容片段:', data.content);
          // 实时显示内容
          break;
        case 'complete':
          console.log('处理完成:', data.metadata);
          break;
        case 'error':
          console.error('错误:', data.error);
          break;
      }
    }
  }
}
\end{lstlisting}

\section{数据模型}

\subsection{会话模型（Session）}

\begin{lstlisting}[language=json]
{
  "id": "string",           // 会话ID
  "title": "string",        // 会话标题
  "preview": "string",      // 会话预览
  "updatedAt": "string"     // 更新时间（ISO 8601）
}
\end{lstlisting}

\subsection{消息模型（Message）}

\begin{lstlisting}[language=json]
{
  "id": "string",           // 消息ID
  "role": "string",         // 角色：user/assistant/system/tool
  "content": "string",      // 消息内容
  "createdAt": "string",    // 创建时间（ISO 8601）
  "status": "string",       // 状态：complete/error
  "metadata": {             // 元数据（可选）
    "tool": "string",       // 使用的工具名称
    "args": {},             // 工具参数
    "iterations": 0         // 工具调用迭代次数
  }
}
\end{lstlisting}

\section{使用示例}

\subsection{完整流程示例}

以下示例展示如何创建一个会话并进行对话：

\begin{lstlisting}[language=javascript]
// 1. 创建会话
const createResponse = await fetch('/api/agent/sessions', {
  method: 'POST',
  headers: { 'Content-Type': 'application/json' },
  body: JSON.stringify({ title: '安全报告生成' })
});
const session = await createResponse.json();
console.log('会话ID:', session.id);

// 2. 发送消息
const chatResponse = await fetch('/api/agent/chat', {
  method: 'POST',
  headers: { 'Content-Type': 'application/json' },
  body: JSON.stringify({
    sessionId: session.id,
    content: '请生成张三的安全报告'
  })
});
const message = await chatResponse.json();
console.log('AI回复:', message.content);

// 3. 获取会话详情
const detailResponse = await fetch(`/api/agent/sessions/${session.id}`);
const detail = await detailResponse.json();
console.log('会话消息数:', detail.messages.length);

// 4. 删除会话
await fetch(`/api/agent/sessions/${session.id}`, {
  method: 'DELETE'
});
\end{lstlisting}

\subsection{流式对话示例}

\begin{lstlisting}[language=javascript]
async function streamChat(content, messages = []) {
  const response = await fetch('/api/agent/chat/stream', {
    method: 'POST',
    headers: { 'Content-Type': 'application/json' },
    body: JSON.stringify({ content, messages })
  });

  const reader = response.body.getReader();
  const decoder = new TextDecoder();
  let fullContent = '';

  while (true) {
    const { done, value } = await reader.read();
    if (done) break;

    const chunk = decoder.decode(value);
    const lines = chunk.split('\n');

    for (const line of lines) {
      if (line.startsWith('data: ')) {
        const data = JSON.parse(line.slice(6));

        if (data.type === 'chunk') {
          fullContent += data.content;
          // 实时更新UI
          updateUI(fullContent);
        } else if (data.type === 'complete') {
          console.log('完成，消息ID:', data.messageId);
        } else if (data.type === 'error') {
          console.error('错误:', data.error);
        }
      }
    }
  }

  return fullContent;
}
\end{lstlisting}

\section{错误处理}

\subsection{常见错误}

\begin{table}[h]
  \centering
  \begin{tabular}{lll}
    \toprule
    错误码 & 错误信息 & 解决方案 \\
    \midrule
    400 & 缺少 sessionId 或 content & 检查请求参数 \\
    404 & 会话不存在 & 确认会话ID是否正确 \\
    500 & 服务器内部错误 & 联系技术支持 \\
    \bottomrule
  \end{tabular}
  \caption{常见错误及解决方案}
\end{table}

\subsection{错误响应示例}

\begin{lstlisting}[language=json]
{
  "error": "会话不存在"
}
\end{lstlisting}

\section{最佳实践}

\subsection{上下文管理}

\begin{itemize}
  \item Agent服务采用无状态设计，不提供会话管理功能
  \item 调用端需要自行管理会话和消息历史
  \item 每次请求时传递完整的历史消息列表（messages参数）
  \item 建议限制历史消息数量，避免超出上下文窗口（默认保留最近50条消息）
  \item 历史消息格式：每个消息包含role（user/assistant）和content字段
\end{itemize}

\subsection{消息处理}

\begin{itemize}
  \item 对于需要实时反馈的场景，使用流式接口（\texttt{/chat/stream}）
  \item 对于批量处理或不需要实时反馈的场景，使用非流式接口（\texttt{/chat}）
  \item 合理控制消息长度，避免超出上下文窗口限制
\end{itemize}

\subsection{错误处理}

\begin{itemize}
  \item 始终检查HTTP状态码
  \item 实现重试机制，处理网络错误
  \item 记录错误日志，便于问题排查
\end{itemize}

\section{附录}

\subsection{工具说明}

\subsubsection{generate\_driver\_report}

生成驾驶员安全评估报告，包含风险评分、行为画像及改进建议。

\textbf{参数}

\begin{itemize}
  \item \texttt{driver\_name} (String, 必填): 驾驶员姓名
\end{itemize}

\subsubsection{generate\_vehicle\_report}

生成车辆运行分析报告，分析车辆的维修记录、故障率及能耗数据。

\textbf{参数}

\begin{itemize}
  \item \texttt{vehicle\_id} (String, 必填): 车辆ID或车牌号
\end{itemize}

\subsubsection{generate\_route\_report}

生成线路风险扫描报告，基于线路的历史事故数据和路况信息。

\textbf{参数}

\begin{itemize}
  \item \texttt{route\_name} (String, 必填): 线路名称或编号
\end{itemize}

\subsubsection{consult\_omni}

智能咨询助手，支持趋势分析、归因分析、政策问答等复杂查询。

\textbf{参数}

\begin{itemize}
  \item \texttt{query} (String, 必填): 用户查询内容
\end{itemize}

\subsection{变更日志}

\begin{table}[h]
  \centering
  \begin{tabular}{lll}
    \toprule
    版本 & 日期 & 变更说明 \\
    \midrule
    1.0.0 & 2024-01-15 & 初始版本 \\
    \bottomrule
  \end{tabular}
  \caption{API版本变更日志}
\end{table}

\end{document}

